\chapter{Architektura systemu operacyjnego i~podział zasobów}
\label{ch:architektura}

Głównym założeniem architektonicznym projektu było rozdzielenie zadań pomiędzy dostępnymi rdzeniami procesora. 

\section{Wydzielenie rdzenia z~systemu Linux (Core Isolation)}
Aby zapewnić przewidywalny czas wykonywania dla kodu bare-metal, system Linux musiał zostać pozbawiony kontroli nad jednym z~rdzeni CPU. Zrealizowano to za pomocą argumentów linii komend jądra (np. parametr \texttt{isolcpus}), co uchroniło dany rdzeń przed standardowym planistą (ang. \textit{scheduler}) systemu operacyjnego.

\section{Mechanizmy komunikacji międzyprocesorowej (IPC)}
Komunikacja pomiędzy systemem Linux uruchomionym na głównym klastrze rdzeni, a~kodem bare-metal działającym na rdzeniu wydzielonym, wymagała niezawodnego mechanizmu IPC (ang. \textit{Inter-Processor Communication}). Wykorzystano do tego bloki pamięci współdzielonej (Shared Memory) oraz przerwania sprzętowe typu IPI (ang. \textit{Inter-Processor Interrupt}).
